% We prove theorem about cusp shapes of hyperbolic knots.

We will now apply the results of the previous sections to determine
bounds on the shapes of cusps of knot complements.  In particular, let
$K$ be a hyperbolic knot with a prime, twist reduced diagram, where
prime and twist reduced are defined in Definitions \ref{def:prime} and
\ref{def:twred}.  We will
bound the normalized height of the cusp of $S^3-K$, using Theorem
\ref{thm:main-norm-height}.  We need the following:
\begin{enumerate}
  \item An initial hyperbolic manifold $X$ with $n+1$ cusps.
  \item A cone deformation deforming $X$ to the knot complement
    $S^3-K$.
  \item A lower bound on the tube radius ($R\geq R_1\geq 0.56$).
\end{enumerate}

Given these pieces, we will obtain bounds on
normalized height.
  
%%%%%%%%%%%%%%%%%%%%%%%%%%%%%%%%%%%%%%%%%%%%%%%%%%%%%%%%%%%%%%%%%
\subsection{An initial manifold}
\label{subsec:init}

% Initial manifold for cone deformation.

Begin with a prime, twist reduced diagram of the knot $K$.  At each
twist region of the diagram, add a closed curve encircling two
strands, called a \emph{crossing circle}.  We now have a diagram of a
link, which we call $J$.  The link complement $S^3-J$ is homeomorphic
to the complement of the link $L$ obtained from $J$ by removing pairs
of crossings from each twist region.  See Figure
\ref{fig:knot-decomp}.

\begin{figure}
\begin{center}
\includegraphics{figures/knot-decomp.eps}
\end{center}
\caption{Left to right:  The link $K$.  The link $K$ with crossing
circles added.  The link $L$.}
\label{fig:knot-decomp}
\end{figure}

These links with crossing circles added have been studied by many
people, including Adams \cite{adams:aug}; Lackenby, Agol and Thurston
\cite{lackenby:alt-volume}.  Many of the properties of these links were
addressed in previous papers (see \cite{purcell:volume}, and
\cite{futer-purcell} with Futer).

The links are useful to us in the cone manifold setting because
first, symmetries of the diagram of $L$ allow us to determine the
hyperbolic structure on $S^3-L$ explicitly.  Secondly, $S^3-L$ is
related to the original knot complement via Dehn filling.  In
particular, if $2m_i$ crossings were removed from the $i$-th twist
region to go from the diagram of $J$ to that of $L$, then by
performing a $1/m_i$ Dehn filling on the corresponding $i$-th crossing
circle of $S^3-L$, we re-insert these crossings and put back the
crossing circle, obtaining $S^3-K$ (see e.g.  \cite{rolfsen-book},
Chapter 9).

Thus, starting with a hyperbolic augmented link $L$, we can determine
information including shapes of the cusps.  Next, we can find
conditions which will ensure the Dehn filling above can be obtained by
a hyperbolic Dehn filling via cone deformation.  Hence we will have
the first two pieces needed to bound the normalized height of the cusp
$K$: The initial manifold $M$ with hyperbolic structure will be
$S^3-L$.  It will be connected by cone deformation to $S^3-K$.

\subsubsection{The geometry of $S^3-L$}

We note in \cite{futer-purcell} that if $K$ has at least $2$ twist
regions, then the augmented link complement $S^3-L$ will admit a
complete, finite volume hyperbolic structure.  This can be shown by
work of Adams \cite{adams:aug}, using the geometrization of Haken
manifolds.  It can also be shown more directly using Andreev's
theorem.  The proof using Andreev's theorem gives a packing of circles
related to $S^3-L$ which we will need again in Section
\ref{subsec:2bridge}, so we include it here.

\begin{theorem}
  If $K$ has a prime, twist reduced diagram with at least $2$ twist
  regions, then the augmented link complement $S^3-L$ will admit a
  complete, finite volume hyperbolic structure.
\label{thm:Lhyp}
\end{theorem}

\begin{proof}
  We prove the lemma by finding the hyperbolic structure.  To do so,
  we decompose $S^3-L$ into two ideal polyhedra, as in
  \cite{lackenby:alt-volume} and \cite{futer-purcell}.  We then find a
  hyperbolic structure on the polyhedra which gives a nonsingular
  hyperbolic structure on $S^3-L$ after gluing.

  To find the polyhedra, first remove all remaining single crossings
  at crossing circles.  That is, if the number of crossings in a twist
  region of the diagram of $K$ was odd, then $S^3-L$ will still have a
  single crossing at that twist region.  Remove it.  The new link has
  components lying flat in the projection plane bound together by
  crossing circles.  Call it $L'$.
  
  Slice $S^3-L'$ along the projection plane.  This cuts the manifold
  into two identical pieces, one on either side of the plane.  Each
  2--punctured disk encircled by a crossing circle has been sliced in
  half.  Slice along these remaining halves of disks, opening them up
  into two faces.  This gives an ideal polyhedral decomposition of the
  manifold.  The edges are given by the intersection of the
  2--punctured disks with the projection plane.  On each of the two
  polyhedra, we have one face per planar region of the diagram of
  $L'$, and two triangular faces per 2--punctured disk.  Shade the
  faces arising from 2--punctured disks, giving each polyhedron a
  checkerboard coloring.  The edges meet in 4--valent ideal vertices.
  At each vertex, two white faces are separated by two shaded faces.
  (See \cite{futer-purcell} or \cite{lackenby:alt-volume} for
  pictures.) 
  
  The gluing pattern on the polyhedra is given by following the above
  process in reverse: First, on each polyhedron fold the pairs of
  triangular shaded faces together at their common vertex.  Then glue
  each white face to its corresponding face on the opposite
  polyhedron via the identity map.  This gives back $S^3-L'$.
  
  To obtain $S^3-L$, we use the same two polyhedra, but make the
  following change in gluing.  For crossing circles encircling a
  single crossing of the diagram of $L$, rather than glue shaded
  triangles across their common vertex on a single polyhedron, we
  insert a half--twist by gluing each triangle of one polyhedron to
  the opposite triangle of the other.  See Figure
  \ref{fig:half-twist}.  

\begin{figure}
  \begin{center}
    \includegraphics{figures/half-twist.eps}
  \end{center}
  \caption{Left:  Gluing 2--punctured discs with no crossings.  Right:
    Gluing 2--punctured discs with a single crossing.}
  \label{fig:half-twist}
\end{figure}
  
Now our goal is to find a hyperbolic structure on these two polyhedra.
We show the two polyhedra actually have totally geodesic faces in
$\HH^3$, meeting in dihedral angles of $\pi/2$.  Since faces are glued
in pairs, and edges in 4's, the gluing will give a complete hyperbolic
structure on $S^3-L$.
  
Note any totally geodesic plane in $\HH^3$ extends to give a circle on
the boundary $S^2$ at infinity.  Conversely, given a collection of
circles on $S^2$, we can view those circles as boundaries of planes in
$\HH^3$.  We obtain our polyhedra by finding appropriate circle
packings, then cutting away half--spaces bounded by hemispheres in
$\HH^3$.  Our tool is a corollary of Andreev's theorem noted by
Thurston in \cite{thurston}:
  
\begin{theorem}[Andreev]
  Let $\gamma$ be a triangulation of $S^2$ such that each edge has
  distinct ends and no two vertices are joined by more than one edge.
  Then there is a packing of circles in $S^2$ whose nerve is isotopic
  to $\gamma$.  This circle packing is unique up to M{\"o}bius
  transformation.
\end{theorem}

Recall that the \emph{nerve} of a circle packing is the graph obtained
by adding a vertex for each circle, and an edge connecting two
vertices if and only if the corresponding circles are tangent.

  
We begin by finding a triangulation of $S^2$ associated with one of
the polyhedra described above.  For each white face, take a vertex.
Connect two vertices by an edge if and only if the two corresponding
white faces meet at a vertex of the polyhedron.  This gives a
triangulation $\gamma$ of $S^2$.  Note if we draw vertices of $\gamma$
on top of white faces of the polyhedron, and edges of $\gamma$ through
the vertices of the polyhedron, that each triangle of $\gamma$
circumscribes a shaded face of the polyhedron.

To apply Andreev's theorem, we need to show this triangulation
$\gamma$ satisfies the conditions: each edge has distinct ends, and no
two vertices are joined by more than one edge.  To show these
conditions, we consider the graph dual to $\gamma$, as follows.

Consider the original link diagram.  We view it as a 4--valent graph.
Replace each twist region of the graph by a single edge, as in Figure
\ref{fig:twist-to-edge}.  

\begin{figure}
  \begin{center}
    \includegraphics{figures/twist_to_edge.eps}
  \end{center}
  \caption{Create a trivalent graph by replacing each twist region with
  an edge.}
  \label{fig:twist-to-edge}
\end{figure}

Since each crossing is part of some twist region, this gives a
trivalent graph $\Gamma$.
%% This probably needs proof.
Note the dual graph to $\Gamma$ is isotopic to the triangulation
$\gamma$.
We use properties of the dual $\Gamma$ to show the triangulation
$\gamma$ satisfies the conditions necessary for Andreev's theorem.

\begin{lemma}
  No edge of $\gamma$ has a single vertex at both endpoints.
  \label{lemma:andreev1}
\end{lemma}

\begin{proof}
Consider the graph $\Gamma$.  If its dual $\gamma$ has an edge with one
vertex at both endpoints, then the graph $\Gamma$ must be of the form
in Figure \ref{fig:Andr_double_end} (a).

\begin{figure}
  \begin{center}
    (a)
  \input{figures/Andr_double_end.pstex_t}
  \hspace{.5in}
  (b)
  \input{figures/Andr_twist_end.pstex_t}
  \end{center}
  \caption{(a.) Form of $\Gamma$ if $\gamma$ has an edge
    with a single vertex at both endpoints.  (b.)  Corresponding form
    of original link diagram.  }
  \label{fig:Andr_double_end}
\end{figure}

Note then if the edge $e$ was an original edge of the 4--valent link
diagram, then we have a simple closed curve intersecting the link
diagram transversely in a single point.  This is impossible.

If instead the edge $e$ replaced a twist region in the original
4-valent link diagram, then the original diagram had the form of
Figure \ref{fig:Andr_double_end} (b).  However, note one of the
subdiagrams $D$ and $D'$ must contain crossings else the original
diagram has at most one twist region.  Without loss of generality, say
$D$ contains crossings.  Then the dashed simple closed curve enclosing
$D$ intersects the link diagram transversely in two points, but both
its interior and its exterior contain crossings.  This contradicts the
fact that the original diagram was prime.
\end{proof}

\begin{lemma}
  No two vertices of $\gamma$ are joined by more than one edge. 
  \label{lemma:andreev2}
\end{lemma}

\begin{proof}
Suppose two vertices of $\gamma$ are joined by two distinct edges.
Then consider the dual graph $\Gamma$.  $\Gamma$ must be of the form
shown in Figure \ref{fig:Andr_double_edge} (a), where region $B$ and
region $B'$ are the same.

\begin{figure}
  \begin{center}
    (a)
  \input{figures/Andr_double_edge.pstex_t}
  \hspace{.5in}
  (b)
  \input{figures/Andr_double_graph.pstex_t}
 
  \end{center}
  \caption{Form of $\Gamma$ when $\gamma$ has two vertices
  connected by two distinct edges.  In (a), the edges of $\gamma$ are
  the dashed lines, with one of the vertices in $A$ and the other at
  $\infty$.}
  \label{fig:Andr_double_edge}
\end{figure}

Give the edge $e_1$ an orientation, say $e_1$ points in the direction
of the edge $a$.  This gives the region $A$ an orientation, and hence
assigns to the edge $e_2$ an orientation.  Now, trace a path of edges
around the region $B$, listing the edges in the order they occur.  Our
list begins with $e_1$, then $a$, etc.  Note $e_2$ must appear in the
list before the edge $d$, else region $B$ would be cut off from region
$B'$.  Also, edge $b$ must occur before edge $c$, by the orientation
on $e_2$.  Hence the graph $\Gamma$ has the form illustrated in Figure
\ref{fig:Andr_double_edge} (b).

Case 1: Edges $e_1$ and $e_2$ both were edges of the original
4--valent link diagram.  Then there is a simple closed curve
intersecting $e_1$ and $e_2$ transversely in the original link diagram
and intersecting the diagram in no other points.  Since the
original diagram has crossings on either side of the region $A$, this
contradicts the fact that we started with a prime diagram.

Case 2: Edges $e_1$ and $e_2$ both replaced twist regions in the
original link diagram.  Then the original link diagram must have been
of the form in Figure \ref{fig:Andr_double_diag1} (a).  Then since the
original diagram was twist reduced, either the subdiagram $D$ or $D'$
must be a string of bigons.  But then rather than two distinct twist
regions, our diagram must have had only one twist region.  Thus edges
$e_1$ and $e_2$ would not have been distinct.

\begin{figure}
  \begin{center}
    (a) \input{figures/Andr_double_diag1.pstex_t}
  \hspace{.5in}
  (b)
  \input{figures/Andr_double_diag2.pstex_t}
  \end{center}
  \caption{Possible forms of the original link diagram if $\gamma$
  contains two vertices connected by two distinct edges.}
  \label{fig:Andr_double_diag1}
\end{figure}


Case 3: One of edge $e_1$ or $e_2$ is from the original link diagram,
the other replaced a twist region.  Then the original link diagram
must have been of the form shown in Figure \ref{fig:Andr_double_diag1}
(b).  Now note that the simple closed curve bounding subdiagram $D$
intersects the link transversely in only three points.  This is
impossible.

Since these are the only three cases, these contradictions prove the
lemma. 
\end{proof}


Now Lemmas \ref{lemma:andreev1} and \ref{lemma:andreev2} are enough to
show that Andreev's theorem applies, and there is a circle packing of
$S^2$ whose nerve is $\gamma$.

Complete the proof of Theorem \ref{thm:Lhyp} by slicing off
half--spaces bounded by geodesic hemispheres in $\mathbb H^3$
corresponding to each circle in the circle packing.  These give the
geodesic white faces of the polyhedron.  The shaded faces are obtained
by slicing off hemispheres in $\mathbb H^3$ corresponding to each
circle of the dual circle packing.  Note white and shaded hemispheres
meet at right angles.  

Thus the polyhedron for the flat augmented link is totally geodesic in
$\mathbb H^3$, and hence the flat augmented link is hyperbolic.
\end{proof}



\subsubsection{Normalized lengths on cusps}  To apply results on cone
deformations, we will need the following information on normalized
lengths of particular curves on the cusps of $S^3-L$.


\begin{prop}
  Let $K$ and $L$ be as above.  In the prime, twist reduced diagram of
  $K$, let $n$ be the number of twist regions.  Let $c_i$ be the
  number of crossings in the $i$-th twist region.  Then on
  cusps of $L$, we have the following normalized lengths.
  \begin{enumerate}
  \item For a cusp corresponding to a crossing circle, let $s_i$ be
    the slope such that {D}ehn filling $S^3-L$ along $s_i$ re-inserts
    the $c_i$ crossings at that twist region.  Then the normalized
    length of $s_i$ is at least $\sqrt{c_i}$.
  \item For the cusp corresponding to the link component in the
    projection plane, the normalized height (i.e. the normalized
    length of the curve orthogonal to a meridian) is at least
    $\sqrt{n}$, and at most $\sqrt{n(n-1)}$.
  \end{enumerate}
\label{prop:proj-height}
\end{prop}

\begin{proof}
  This can be deduced from results of \cite{futer-purcell}.  However,
  given the circle packing associated to $S^3-L$ developed in the
  previous section, we are able to streamline the proof somewhat.
  
  First, recall the polyhedra obtained in the proof of Theorem
  \ref{thm:Lhyp}.  Note each cusp will be tiled by rectangles given by
  the intersection of the cusp with the totally geodesic white and
  shaded faces of the polyhedra.  Two opposite sides of each of these
  rectangles come from the intersection of the cusp with shaded faces
  of the polyhedra, or from the 2--punctured disks in the diagram of
  $L$, and the other two sides come from white faces.  Call the sides
  shaded sides and white sides, respectively.  
  
  For part (1), note any crossing circle intersects a single
  2--punctured disk in a longitude.  Half of the longitude is given by
  the intersection with one polyhedron, the other half by intersection
  with the other.  The crossing circle intersects no other shaded
  faces.  Thus the crossing circle cusp is tiled by exactly two
  identical rectangles.
  
  When the crossing circle encircles no single crossing, or half twist,
  then two white faces intersect the cusp in meridians.  In this case,
  $c_i$ is even, $c_i=2m_i$, and the slope $s_i$ is $1/m_i$.  It is
  given by one step along a white side of a rectangle (a meridian),
  plus $2m_i=c_i$ steps along shaded sides of the rectangle.
  
  When the crossing circle encircles a single crossing, then the two
  rectangles tiling the cusp are sheared.  In this case, a meridian is
  given by one step along a white side, plus (or minus) a step along a
  shaded side.  Now $c_i =2m_i+1$ is odd, and the slope $s_i =1/m_i$
  is given by a step along a white side, plus (or minus) a step along
  a shaded side, plus (or minus) $2m_i$ steps along shaded sides.
  
  In either case, the slope is given by a single step along a white
  side plus (or minus) $c_i$ steps along a shaded side.  Let $w$
  denote the length of a white side, and let $s$ denote the length of
  a shaded side.  Thus the normalized length of $s_i$ is:
  $$\frac{\sqrt{w^2+c_i^2s^2}}{\sqrt{2ws}} =
  \sqrt{\frac{w}{2s}+c_i^2\frac{s}{2w}}.$$
  Let $u=w/(2s)$.  Then for
  constant $c_i$, this quantity is minimized when $u=c_i/2$, and the
  minimum value is $\sqrt{c_i}$.
  
  For part (2), we need to know more about the possible lengths of the
  white sides of rectangles tiling the cusp.  For a given rectangle,
  consider the corresponding vertex in one of the polyhedra.  On the
  sphere $S^2$ at infinity, this vertex will look like a point of
  tangency of two circles in the packing given by Andreev's theorem.
  Since the nerve of this circle packing is a triangulation of $S^2$,
  the point of tangency corresponds to an edge on two distinct
  triangles of this nerve.  Thus there are two additional circles
  tangent to the two given circles, as in Figure \ref{fig:circles}.

\begin{figure}
  \begin{center}
    \includegraphics{figures/circles1.eps} \hspace{.2in}
    \includegraphics{figures/circles2.eps}
  \end{center}
  \caption{The form of the packing of circles about any vertex.
    Dashed circles correspond to the shaded faces of the polyhedron at
    this vertex.}
  \label{fig:circles}
\end{figure}

Now, apply a M{\"o}bius transformation taking the vertex, or the point
of tangency of the circle packing, to infinity.  This takes the two
tangent circles to two parallel lines.  It takes the two additional
tangent circles to circles tangent to both the parallel lines, as in
Figure \ref{fig:rec-circles}.  Note this also gives the similarity
structure of the rectangle under consideration.  If we normalize so
that the shaded side (coming from the intersection with a 2--punctured
disk) has length $1$, then the circle lying under the dashed line in
Figure \ref{fig:rec-circles} has diameter $1$.  Since circles in the
circle packing do not overlap, this forces a white side (i.e. without
dashes in the figure) to have length at least $1$.

\begin{figure}
  \begin{center}
    \includegraphics{figures/rec-circ1.eps}
    \hspace{.2in}
    \includegraphics{figures/rec-circ2.eps}
  \end{center}
  \caption{Result of taking point of tangency to infinity.}
  \label{fig:rec-circles}
\end{figure}

For an upper bound on the length of a white side, note it is longest
when all circles in the circle packing are lined up in a row, tangent
to the two parallel lines.  In this case, when a shaded side has
length $1$, the length of the white side will be the number of circles
between the parallel lines, minus $1$.  To find the number of circles
possible, we use an Euler characteristic argument.  Note there will be
$2n$ shaded faces in one of the polyhedra, coming from opening up $n$
crossing circles.  These each correspond to a single face of the
triangulation of $S^2$ used in Andreev's theorem.  Because it is a
triangulation, the number of edges equals $3/2$ the number of faces,
or $3n$.  Then by an Euler characteristic argument, there are $n+2$
vertices.  Recall each vertex gave a circle in our circle packing.
Since two of these circles become our parallel lines, there are $n$
circles packed between these.  Thus the maximum length of a white side
is $n-1$.

Now, note that the cusp in the projection plane intersects each
2--punctured disk $2n$ times, each time in a meridian.  Half of this
meridian will come from the intersection of one polyhedron (the
``top'' half), the other half from the intersection of the other
polyhedron.  Thus the meridian runs along two shaded sides of
rectangles.  Still normalizing so each shaded side has length $1$,
this implies a meridian has length $2$.  The curve giving the height
$H$ of the cusp runs along $2n$ white sides of the rectangles (one for
each 2--punctured disk of intersection).  The normalized height will
be given by $H/\sqrt{2H} = \sqrt{H/2}.$ The height $H$ will be minimal
when each of the $2n$ white sides it runs along are of minimal length,
or length $1$.  Thus the minimal normalized height is $\sqrt{n}$.

Similarly, the height $H$ will be maximal when each white side is of
maximal length, or length $n-1$.  Thus the maximal normalized height
is $\sqrt{n(n-1)}$.
\end{proof}


\subsubsection{Hyperbolic cone deformation on $S^3-L$}
  \label{subsec:crossings}

If the normalized lengths of the slopes on which {D}ehn filling was
performed were sufficiently large, then Hodgson and Kerckhoff showed
that the {D}ehn filling could be realized by a hyperbolic cone
deformation \cite{hk:univ}.  Proposition \ref{prop:proj-height}
indicates that to attain these minimal normalized slope lengths, we
need only ensure that there are sufficiently many crossings in each
twist region.  In particular, by \cite{hk:univ}, provided $\sqrt{c_j}
\geq 10.6273$, or provided that there are at least $113$ crossings per
twist region of the prime, twist reduced diagram of $K$, then the
{D}ehn filling of $S^3-L$ can be realized by a hyperbolic cone
deformation.

Actually, we will see that the number of crossings in each twist
region needs to be larger than $113$ to apply Theorem
\ref{thm:bound-bdry-time}.  For now, however, note that we may ensure
a hyperbolic cone deformation does exist from $S^3-L$ to $S^3-K$ by
putting conditions on the diagram of $K$.


%%%%%%%%%%%%%%%%%%%%%%%%%%%%%%%%%%%%%%%%%%%%%%%%%%%%%%%%%%%%%%%%%
\subsection{Finding values for tube radius}
\label{subsec:tube}

We need to ensure that the tube radius remains larger than $R_1 =
0.56$ throughout the cone deformation.  We can do so by again applying
work of Hodgson and Kerckhoff \cite{hk:univ}.  They show that in the
neighborhood of the $j$-th component of the singular locus, we can
increase the cone angle from $0$ to $2\pi$, maintaining the tube
radius $R \geq R_1$, provided equation (\ref{eq:Lhat_int}) of the
proof of Lemma \ref{lemma:b_j-bound} holds.

That is, consider a particular component of the singular locus
$\Sigma_j$.  The singular locus consists of crossing circles of the
link $L$.  At the $j$-th crossing circle, the hyperbolic cone
deformation performs {D}ehn filling along a slope of the form $1/n_j$.
Let $L_0$ denote the normalized length of this slope in $S^3-L$.  We
can then increase the cone angle from $0$ to $2\pi$, maintaining the
tube radius bound $R \geq R_1$, provided the normalized length $L_0^2$
is larger than the increasing function $I(R_1)$ defined in equation
(\ref{eq:Lhat_int}).

So to ensure $R \geq R_1 \geq 0.56$, we compute that $L_0^2$ must be
at least $113.044$.  By Proposition \ref{prop:proj-height}, we know
$L_0^2$ is at least $c_i$, where $c_i$ is the number of crossings in
the $i$-th twist region.  Thus by ensuring there are at least $114$
crossings in the $i$-th twist region, we will guarantee that our tube
radius is at least $0.56$.  For general $R_1$, we have the following
result.

\begin{lemma}
  Let $K$, $L$ be as above, such that we have a cone deformation from
  $S^3-L$ to $S^3-K$.  Fix $R_1 \geq 0.531$.
  Then $R \geq R_1$ throughout the deformation provided the number of
  crossings in each twist region is at least $I(R_1)$, where $I(R_1)$
  is defined in equation (\ref{eq:Lhat_int}).
\label{lemma:crossings-vs-R}
\end{lemma}


%%%%%%%%%%%%%%%%%%%%%%%%%%%%%%%%%%%%%%%%%%%%%%%%%%%%%%%%%%%%%%%%%
\subsection{Proof of Theorem \ref{thm:main-knot-cusp}}
\label{subsec:numeric}

\hspace{.1in}

\medskip

\noindent\textbf{Theorem \ref{thm:main-knot-cusp}} \emph{
  Let $K$ be a knot in $S^3$ which admits a prime, twist reduced
  diagram with a total of $n\geq 2$ twist regions, with each twist
  region containing at least $c\geq 116$ crossings.  In a hyperbolic
  structure on $S^3-K$, take a horoball neighborhood $U$ about $K$.
  Normalize so that the meridian on $\partial U$ has length $1$.  Then
  the height $H$ of the cusp of $S^3-K$ satisfies:
  $$H\geq n\,(1-f(c))^2.$$
  Here $f(c)$ is a positive function of $c$ which
  approaches $0$ as $c$ increases to infinity.
Additionally, $H\leq n(\sqrt{n-1} + f(c))^2$.
}


\begin{proof}
  
  If $K$ admits a prime, twist reduced diagram with at least $2$ twist
  regions, then we may form the augmented link $L$ from the diagram of
  $K$.  We have seen in \S\ref{subsec:crossings} that, provided there
  are at least $113$ crossings in each twist region of the diagram of
  $K$, then $S^3-L$ and $S^3-K$ are connected by a hyperbolic cone
  deformation $X_t$, with $X_0=S^3-L$, $X_\tau=S^3-K$.  The singular
  locus will consist of the crossing circle components of $L$.  There
  will be $n$ of these.
  
  By results of Section \ref{subsec:tube}, if there are at least $114$
  crossings per twist region then the tube radius of the deformation
  is at least $R_1=0.56$.  Then by Theorem \ref{thm:main-norm-height}, 
  $$ -(2\pi)^2\frac{\sqrt{n\,C(R_1)}}{(1-e^{-2R_1})\sqrt{2}}
  \leq h_\tau-h_0
  \leq
  (2\pi)^2\frac{\sqrt{n\,C(R_1)}}{(1-e^{-2R_1})\sqrt{2}},$$
  where $h_t$ denotes the normalized height of the cusp at time $t$.
  
And so:
\begin{equation}
 h_0 - (2\pi)^2 \frac{\sqrt{n\,C(R_1)}}{(1-e^{-2R_1})\sqrt{2}} \leq
h_\tau \leq
h_0 + (2\pi)^2 \frac{\sqrt{n\,C(R_1)}}{(1-e^{-2R_1})\sqrt{2}}.
\label{eqn:2-sides}
\end{equation}

We may now apply part (2) of Proposition \ref{prop:proj-height}.  This
implies that the initial normalized height $h_0$, i.e. the normalized
height on $S^3-L$, is between $\sqrt{n}$ and $\sqrt{n(n-1)}$.

Plug this into the left hand side of the above equation and we find:
\begin{equation}
h_\tau \geq \sqrt{n}\left(1-(2\pi)^2
  \frac{\sqrt{C(R_1)}}{(1-e^{-2R_1})\sqrt{2}}\right).
\label{eqn:h}
\end{equation}

For the right hand side, we find:
\begin{equation}
  h_\tau \leq \sqrt{n}\left(\sqrt{n-1} + (2\pi)^2
  \frac{\sqrt{C(R_1)}}{(1-e^{-2R_1})\sqrt{2}}\right).
\label{eqn:hright}
\end{equation}

The term $C(R_1)$ can be computed explicitly using the formulas in the
proof of 
Lemma \ref{lemma:b_j-bound}.  Recall $C(R_1)$ is strictly decreasing,
approaching $0$ as $R_1$ approaches infinity.  Thus the function
$$\bar{f}(R_1)= (2\pi)^2\frac{\sqrt{C(R_1)}}{(1-e^{-2R_1})\sqrt{2}},$$
which appears in both (\ref{eqn:h}) and (\ref{eqn:hright}), is
decreasing with $R_1$, approaching $0$.  The graph of $\bar{f}$, as
well as that of $I(R)$, is shown in Figure \ref{fig:F-I-graphs}.

\begin{figure}
  \begin{center}
    \framebox{
      \begin{tabular}{c}
    \includegraphics[height=1.6in]{figures/F-R1.eps}\\
    $y=\bar{f}(R)$
    \end{tabular}}
\hspace{.1in}
\framebox{
  \begin{tabular}{c}
    \includegraphics[height=1.6in]{figures/I-R1.eps}\\
    $y=I(R)$
  \end{tabular}}
  \end{center}
  \caption{Graphs of $\bar{f}(R)$ and $I(R)$.  Here $\bar{f}(R)$ gives the
    bounds in Theorem \ref{thm:main-knot-cusp} for given tube radius
    $R$.  $I(R)$ gives the number of crossings required to guarantee
    tube radius $R$.}
  \label{fig:F-I-graphs}
\end{figure}

When $R_1=0.56$, $1-\bar{f}(R_1)$ is negative.  Since $h_\tau$
is known to be positive, this doesn't give a very good estimate in
(\ref{eqn:h}).  To get a better estimate, we increase $R_1$.  The
value $1-\bar{f}(R_1)$ becomes positive when $R_1$ is about $0.6624$.
However, as $R_1$ increases, so does $I(R_1)$, so by Lemma
\ref{lemma:crossings-vs-R}, we need to increase the number of
crossings in each twist region to guarantee that the tube radius
remains larger than $R_1$ throughout the deformation.  In particular,
$c$ must be at least $116$ to ensure $R_1 \geq 0.6624$.

In general, let $f(c) = \bar{f}(I^{-1}(c))$.  Since $I(R)$ is
strictly increasing for $R\geq 0.56$, this is well defined.  Then for
$c\geq 116$, $1-f(c)$ is positive.

To finish the proof, note that we can scale so that the meridian of
the cusp of $S^3-K$ has length $1$.  Then the actual height $H$ of the
cusp of $S^3-K$ will be at least $(h_\tau)^2$.
\end{proof}

Notice in Figure \ref{fig:F-I-graphs} that $\bar{f}$ is close to $0$
for values of $R$ near $2$.  However for $R$ this size, $I(R)$, or the
number of crossings required per twist region, is nearly $700$.  We
find a nice compromise by choosing $R=1.0$.  Then $I(R)$ is not much
larger than $116$, and yet $1-\bar{f}$ is already over $0.8$.  We use
this fact in the following corollary, which gives more explicit
numbers to the results of Theorem \ref{thm:main-knot-cusp}.

\begin{corollary}
  Let $K$ be as in the previous theorem, with at least $c\geq 145$
  crossings in each twist region.  Then when the meridian of the cusp
  of $S^3-K$ is normalized to have length $1$, the height of the cusp
  satisfies: $$H \geq n\,(0.81544)^2.$$
\label{cor:0.8}
\end{corollary}

\begin{proof}
  Letting $c\geq 145$ in the above theorem ensures that $R> R_1 =
  1.0$.  Then $1-f(c) = 1-\bar{f}(1.0) \geq 0.81544$.
\end{proof}



%%%%%%%%%%%%%%%%%%%%%%%%%%%%%%%%%%%%%%%%%%%%%%%%%%%%%%%%%%%%%%%%%
\subsection{Example: 2--bridge knots}
\label{subsec:2bridge}

The requirement of Theorem \ref{thm:main-knot-cusp} that there be at
least $116$ crossings per twist region was based on the worst case
analysis.  This came from Proposition \ref{prop:proj-height}, in which
we showed that in the worst possible case, the normalized length of a
slope was at least the square root of the number of crossings.  When
we are dealing with particular classes of knots, we can often make
significantly better estimates, and reduce the number of crossings
required.  In this section, we present a particular example of such a
class of knots:  2--bridge knots.

We will determine the circle packing obtained from the proof of
Theorem \ref{thm:Lhyp} for a 2--bridge knot.  A 2--bridge knot with an
even number of twist regions will have a corresponding augmented link
$L$ of the form of Figure \ref{fig:2bridge}, except possibly with
single crossings added at some crossing circles.  The picture for an
odd number of twist regions is similar.

\begin{figure}
  \begin{center}
  \input{figures/2bridge.pstex_t}
  \end{center}
  \caption{Augmented link of a 2--bridge knot with an even number
    of twist regions, and an even number of crossings per region.}
  \label{fig:2bridge}
\end{figure}

In any case, the link will correspond to a trivalent graph $\Gamma$ of
the form shown in Figure \ref{fig:2bridge_tri}.  Note each region of
the graph complement gives a circle in the circle packing, with
tangencies across edges. Thus each circle will be tangent to the
circles corresponding to regions $A$ and $D$.  The circle packing
associated with this graph is of the form in Figure
\ref{fig:2bridge_circ}.

\begin{figure}
  \begin{center}
  \input{figures/2bridge_tri.pstex_t}
  \end{center}
  \caption{Trivalent graph associated with a 2--bridge knot.}
  \label{fig:2bridge_tri}
\end{figure}

\begin{figure}
  \begin{center}
  \input{figures/Fig_2bridge_circ.pstex_t}
  \end{center}
  \caption{Circle packing associated with a 2--bridge knot.}
  \label{fig:2bridge_circ}
\end{figure}

Recall that the shape of cusps can be determined by analyzing the
rectangles we get from taking points of tangency of the circles in
Figure \ref{fig:2bridge_circ} to infinity.  In particular, we are
interested in the cusp shapes which correspond to crossing circles in
the link $L$.  The rectangles tiling these cusps come from the
tangencies of circle $A$ with certain of the $B_i$ and $C_i$, and from
tangencies of $D$ with certain of the $B_i$ and $C_i$.  Thus we find
the general shape of rectangles of this form.

When we take a point of tangency between $A$ and $B_i$ to infinity,
$i=1$ or $2$, the two circles tangent to both $A$ and $B_i$ are $D$
and a $C_j$.  These are tangent to each other.  Thus our rectangle in
this case is a square.  Similarly for points of tangency between $D$
and $B_i$.

A crossing circle coming from one of these points of tangency will
therefore be tiled by two squares, with a longitude running across two
shaded sides.  Then the slope $s_i$ along which Dehn filling is
performed will have normalized length $\sqrt{1+c_i^2}/\sqrt{2}$.  Note
this is larger than $I(0.6624)$, implying that the results of Theorem
\ref{thm:main-knot-cusp} apply, provided $c_i \geq 16$.  It is larger
than $I(1.0)$, implying the results of Corollary \ref{cor:0.8} hold,
when $c_i \geq 17$.  Thus the number of crossings in our theorem has
gone from $116$ to $16$, and $145$ to $17$.

When we take a point of tangency between $A$ and $C_i$, or between $D$
and $C_i$ to infinity, we get a $1$ by $2$ rectangle.  This is because
the circles $A$ and $C_i$ are tangent to two distinct circles
$C_{i-1}$ and $C_{i+1}$, or possibly a $B_i$, which are each tangent
to $D$.  Since $D$ is also tangent to $A$ and $C_i$, in the picture at
infinity we have two parallel lines, $A$ and $C_i$, with three full
sized circles tangent to each other and also to the lines $A$ and
$C_i$.

In this case, the crossing circle is tiled by two $1$ by $2$
rectangles, stacked to form a square.  The slope $s_i$
will have normalized length $\sqrt{4+c_i^2}/\sqrt{4}$.  It will be
larger than $I(0.6624)$ provided $c_i \geq 22$, and larger
than $I(1.0)$ provided $c_i \geq 24$.

%% Add info on initial normalized height
We can also make improvements on the initial normalized height,
i.e. the value $h(X)$ of Theorem \ref{thm:main-norm-height}.  To
compute $h(X)$, we determine the shape of the rectangles tiling the
cusp corresponding to the knot strand.  These include all rectangles
obtained by taking a point of tangency of the circle packing of Figure
\ref{fig:2bridge_circ} to infinity, except those tangencies
corresponding to crossing circle cusps.

Each circle tangency of the form $C_{2k}$ tangent to $D$ will give a
rectangle on the knot strand cusp.  Similarly, those tangencies of
circles $C_{2k+1}$ tangent to $A$ will give rectangles on the knot
strand cusp.  As we saw above, these rectangles will each have
dimensions $1$ by $2$.  The side of length $2$ will contribute to the
height.  There will be $(n-2)$ of these rectangles, since each $C_j$
contributes one.  Thus these rectangles contribute $2n-4$ to the
height.

Each circle tangency $C_i$ tangent to $C_{i+1}$ also gives a rectangle
on the knot strand cusp, as do the tangencies $B_1$ with $C_1$ and
$C_{n-2}$ with $B_2$.  Each of these rectangles is a square,
contributing $1$ to the height, and there are $(n-1)$ such rectangles.

The tangencies of $B_1$ and $D$ and of $B_2$ and $A$ (or $D$ if $n$ is
odd) also contribute rectangles of dimension $1$ by $1$ to the knot
strand.  Finally, the tangency of $A$ to $D$ gives a rectangle with
dimensions $1$ by $(n-1)$, contributing $(n-1)$ to the height.  Thus
the total height of the cusp is
$$(2n-4) + (n-1) + 2 + (n-1) = 4n -4.$$

(When $n=2$ this argument needs to be modified slightly:  there are no
$C_i$'s, but $B_1$ and $B_2$ are tangent.  Tangencies of $A$ and
$B_2$, $D$ and $B_1$, and $A$ and $D$ also contribute to the height.
These all give $1$ by $1$ rectangles, hence the height is $4 = 4n-4$
in this case as well.)

A meridian runs along the widths of two rectangles in the cusp
tiling.  Hence it has length $2$.  Thus the normalized height $h(X)$
is given by:
$$h(X) = \frac{4n-4}{\sqrt{(4n-4)(2)}} = \sqrt{2(n-1)}.$$

This information will show:


\begin{prop}
  Let $K$ be a 2--bridge knot in $S^3$ which admits a prime, twist
  reduced diagram with $n\geq 2$ twist regions and at least $c\geq 24$
  crossings in each twist region.  In a hyperbolic structure on
  $S^3-K$, normalized such that the meridian has length $1$, the
  height of the cusp of $S^3-K$ satisfies:
  $$\left(\sqrt{2(n-1)} - \sqrt{n}(0.18456)\right)^2 \leq H \leq
	\left(\sqrt{2(n-1)} + \sqrt{n}(0.18456)\right)^2.$$
\label{prop:2bridge}
\end{prop}

\begin{proof}
Let $L$ be the augmented link corresponding to $K$.  Provided we have
at least $n\geq 2$ twist regions, Theorem \ref{thm:Lhyp} implies
$S^3-L$ is hyperbolic.  We have seen in the above discussion that,
provided we have at least $c\geq 24$ crossings in each twist region,
the slope on each crossing circle of $L$ on which we perform Dehn
filling has normalized length at least $I(1.0)$.  This implies that a
cone deformation exists from $S^3-L$ to $S^3-K$ for which the tube
radius is at least $R_1= 1.0$ throughout.  Then Theorem
\ref{thm:main-norm-height} implies the normalized height of the cusp
of $S^3-K$ satisfies:
$$-\frac{\sqrt{n}(2\pi)^2\sqrt{C(1.0)}}{(1-e^{-2.0})} \leq
 h_{\tau} - \sqrt{2(n-1)} \leq
 \frac{\sqrt{n}(2\pi)^2\sqrt{C(1.0)}}{(1-e^{-2.0})}.$$
Using the formulas for $C(1.0)$ from the proof of Lemma
 \ref{lemma:b_j-bound}, we find
$$\sqrt{2(n-1)} - \sqrt{n}(0.18456) \leq h_{\tau} \leq
 \sqrt{2(n-1)} + \sqrt{n}(0.18456).$$
We obtain the final result by noting if we scale such that the
meridian has length $1$, then the height is at least $(h_\tau)^2$.
\end{proof}





